\scanpage{44}
\mspage{42}
% archive sheet number: 43
pour $\forall\,1\leq i\leq n$, $S'$ est modérément ramifié en les pts \deleted{\unsure{résiduels du diviseur}} \inserted{génériques de $\operatorname{div}(x_i)$}.

d) L'action naturelle de $\pi_1(U,\unsure{\xi|\zeta})$ sur $\inserted{I_t=}\prod_{\ell\ne p}I/[I,I]^{(\ell)}$, qui se factorise \deleted{\unsure{par}} évidemment par une action de $\operatorname{Gal}(\overline k/k)$, s'identifie via l'isomorphisme (8.5.5) à l'action \deleted{de $\operatorname{Gal}(\overline k/k)$} déduite de l'action naturelle de $\operatorname{Gal}(\overline k/k)$ sur $T_\ell(\overline k^{\unsure{\times|*}})$, \deleted{\unsure{déduite}} via les opérations de $\operatorname{Gal}(\overline k/k)$ sur les ${}_{\ell^\nu}(\overline k^*)\subset\overline k^*$.

Démonstration laissée au lecteur \unsure{---}.

\uline{Corollaire} 8.6. Sous les conditions de 8.5, supposons qu'aucune extension finie de $k$ ne contienne toutes les racines $\ell^\nu$-ièmes de $1$, pour $\forall$ entier $\nu$. Alors, pour $\forall$ représentation $\ell$-adique de $\pi_1(U,\xi)$ \inserted{``modérée'', i.e. telle que} l'action induite de $I$ \deleted{est quasi-unipotente} \inserted{se factorise par une action de $I_t$, cette dernière est} \uline{quasi-unipotente}.

\uline{Remarque 8.7} \unsure{a)} Dans le cas d'application de 8.6, on aura
\deleted{$t=u\prod x_i^{\alpha_i}$, avec des $\alpha_i\geq1$, et $x_i\in\unsure{\mathfrak m_i|m_i}\subset\mathfrak m=tV$, \inserted{$u_\alpha\in A^*\subset V^*$,} ce qui prouve que $n=1$ (un seul $x_i$ \inserted{soit $x_1$}) et $\alpha_1=1$. Donc \unsure{on a}}
\deleted{Alors la connaissance de la structure du groupe}
$n=1$, et l'hom $I\longrightarrow I'$ induit un \uline{isom} $I_t\xrightarrow{\sim}I'_t$.

\uline{Remarque} \deleted{\unsure{b}} 8.8. La démonstration de 8.1. marche sans hyp. de propreté sur $X_K$, si on sait que l'on peut trouver un voisinage ouvert $U_\alpha$ du pt générique $\eta_\alpha$ de $S_\alpha$
